
\chapter{射影态射与 Spec、Proj 的相对版本}
\label{ch:projective-morphisms}

本章将运用\trUniversalProperty{}来定义两个有用的构造：\trScheme{} $X$ 上\trSheafOfAlgebras（sheaf of algebras）$\mathscr{R}$ 的 $\mathit{Spec}$，以及\trSheafOfGradedAlgebras（sheaf of graded algebras）$\mathscr{S}_\bullet$ 的 $\mathit{Proj}$。它们将把我们对 $A$-代数之 $\Spec$ 与分次 $A$-代数之 $\Proj$ 的既有构造加以推广（整体化）。我们将看到，\trAffine\trMorphism{}恰是形如 $\mathit{Spec}\mathscr{R}\to X$ 的态射；相应地，我们将把\trProjective\trMorphism{}定义为形如 $\mathit{Proj}\mathscr{S}_\bullet\to X$ 的态射。

在两种情形，我们的计划都是把一个我们在环上熟知的观念，改造成在\trScheme{}上一般地可行。主要的问题在于：如何把各\trAffine{}开子集上的构造粘合起来。我们采取的漂亮路线是：先给出一个\trUniversalProperty；再证明\trAffine{}构造满足这个\trUniversalProperty；然后证明这个\trUniversalProperty{}对\trOpenSubset{}表现良好；最后借助当初让我们得以粘合\trFiberProduct（或\trNormalization）的那个想法，来完成全部艰巨的粘合工作。这个计划里最恼人的部分，是找到正确的\trUniversalProperty——在 $\mathit{Proj}$ 的情形尤其如此。


\section{（拟凝聚）代数层的相对 $\Spec$}

给定一个 $A$-代数 $R$，我们可以取它的 $\Spec$，从而得到 $\Spec A$ 上的一个\trAffineScheme：$\Spec R\to\Spec A$。现在我们将看到这种记号之整体化的一个\trUniversalProperty{}描述。考虑任意一个\trScheme{} $X$，及其上的一个\trQC{}\trSheafOfAlgebras（quasicoherent sheaf of algebras）$\mathscr{R}$。我们将定义如何取这个\trSheafOfAlgebras{}的 $\mathit{Spec}$，并且我们将得到一个「在 $X$ 上仿射」的\trScheme{} $\mathit{Spec}\mathscr{R}\to X$，也就是说，其\trStructureMorphism{}是一个\trAffine\trMorphism。你可以通过两种方式来领会它。

\textbf{17.1.1.}\quad 首先，也是最具体地：对任何\trAffine{}开集 $\Spec A\subset X$，$\Gamma(\Spec A,\mathscr{R})$ 是某个 $A$-代数；称其为 $R$。于是在 $\Spec A$ 之上，$\mathit{Spec}\mathscr{R}$ 将会是 $\Spec R$。

\textbf{17.1.2.}\quad 其次，它将满足一个\trUniversalProperty。我们本可以凭借如下事实来定义 $A$-\trScheme{} $\Spec R$：到 $\Spec R$ 的\trMorphism（从 $\Spec A$ 上的 $A$-\trScheme{} $W$ 出发）与 $A$-代数的映射 $R\to\Gamma(W,\OO_W)$ 相对应（这是我们的老朋友\mbox{练习~7.3.G}）。$\beta\colon\mathit{Spec}\mathscr{R}\to X$ 的\trUniversalProperty{}推广了这一点。给定一个\trMorphism{} $\mu\colon W\to X$，$X$-态射 $W\to\mathit{Spec}\mathscr{R}$ 与使交换图
\begin{equation*}\tag{17.1.2.1}\label{eq:relative-spec-univ-diagram}
\begin{tikzcd}[row sep=1.7em, column sep=2.6em]
& \OO_X \ar[ld] \ar[rd] & \\
\mathscr{R} \ar[rr, "\alpha"] && \mu_*\OO_W
\end{tikzcd}
\end{equation*}
得以交换的诸\trSheafOfAlgebras{}之间的\trMorphism{} $\alpha$ 以（对 $W$）函子的方式一一对应。这里，映射 $\OO_X\to\mu_*\OO_W$ 来自赋环空间映射，而映射 $\OO_X\to\mathscr{R}$ 则来自 $\mathscr{R}$ 的 $\OO_X$-代数结构。（致专家：$\mu_*\OO_W$ 未必是\trQC{}的，但这无关紧要。非专家请完全无视这条括注说明。）

凭一通\trUniversalProperty{}式的废话，这组数据便把 $\beta\colon\mathit{Spec}\mathscr{R}\to X$ 在相差唯一\trIsomorphism{}的意义下确定了下来——前提是它存在。

花哨的翻译：在 $X$-\trScheme{}的\trCategory{}中，$\beta\colon\mathit{Spec}\mathscr{R}\to X$ 表示函子
\[
(\mu\colon W\to X)\longmapsto\{(\alpha\colon\mathscr{R}\to\mu_*\OO_W)\}.
\]

\textbf{17.1.A.（练习）}\quad 证明：若 $X$ 是\trAffine{}的，比如说 $\Spec A$，而 $\mathscr{R}=\widetilde{R}$，其中 $R$ 是一个 $A$-代数，则 $\Spec R\to\Spec A$ 满足这个\trUniversalProperty。（提示：\mbox{练习~7.3.G}。）

\textbf{17.1.3.\ \trhProposition{}。}\ ——\ \emph{设 $\beta\colon\mathit{Spec}\mathscr{R}\to X$ 满足 $(X,\mathscr{R})$ 的\trUniversalProperty，而 $i\colon\mathfrak{U}\hookrightarrow X$ 是一个\trOpenSubscheme。那么 $\beta|_{\mathfrak{U}}\colon\mathit{Spec}\mathscr{R}\times_X\mathfrak{U}=(\mathit{Spec}\mathscr{R})|_{\mathfrak{U}}\to\mathfrak{U}$ 便满足 $(\mathfrak{U},\mathscr{R}|_{\mathfrak{U}})$ 的\trUniversalProperty。}

\emph{证明。}\quad 设 $W$ 是一个 $\mathfrak{U}$-\trScheme，带有\trStructureMorphism{} $\nu\colon W\to\mathfrak{U}$。$\mathfrak{U}$-态射 $W\to\mathit{Spec}\mathscr{R}\times_X\mathfrak{U}$ 与 $X$-态射 $W\to\mathit{Spec}\mathscr{R}$ 是同一回事（这里我们已经由假设知道 $\mu=i\circ\nu\colon W\to X$ 经由 $\mathfrak{U}$ 分解）：
\begin{tikzcd}[row sep=2.3em, column sep=2.6em]
W \ar[rd] \ar[ddr, bend right=28, "\nu"'] & & \\
& \mathit{Spec}\mathscr{R}\times_X\mathfrak{U} \ar[r] \ar[d] & \mathit{Spec}\mathscr{R} \ar[d, "\beta"] \\
& \mathfrak{U} \ar[r, hook, "i"] & X
\end{tikzcd}
由 $\mathit{Spec}\mathscr{R}$ 的\trUniversalProperty，这与一个映射 $\mathscr{R}\to\mu_*\OO_W$ 是同一信息；而由 $(\mu^*,\mu_*)$ 的\trAdjointness（\S14.5.5），它与 $\mu^*\mathscr{R}\to\OO_W$ 是同一回事；又（由 $\nu^*(\mathscr{R}|_{\mathfrak{U}})=\nu^*i^*\mathscr{R}=\mu^*\mathscr{R}$）它就是 $\nu^*(\mathscr{R}|_{\mathfrak{U}})\to\OO_W$。再由 $(\nu^*,\nu_*)$ 的\trAdjointness，这与 $\mathscr{R}|_{\mathfrak{U}}\to\nu_*\OO_W$ 是同一回事，\trUniversalProperty{}得证。\hfill$\square$

把\mbox{练习~17.1.A}与\mbox{命题~17.1.3}结合起来，我们便在如下情形证明了 $\mathit{Spec}\mathscr{R}$ 的存在性：$Y$ 是一个\trAffineScheme{}的\trOpenSubscheme，而 $\mathscr{R}$ 是那个\trAffineScheme{}上某个\trQC{}\trSheafOfAlgebras{}限制到 $Y$ 上的结果。

\textbf{17.1.B.（练习）}\quad 证明 $\mathit{Spec}\mathscr{R}$ 在一般情形下的存在性，其思路循着我们构造\trFiberProduct、\trNormalization{等等}的那套哲学。

\textbf{17.1.4.}\quad 我们来做几条快速的观察。首先，$\mathit{Spec}\mathscr{R}$ 可以「在 $X$ 上\trAffineLocal{}地计算」。我们还有一条\trIsomorphism{} $\phi\colon\mathscr{R}\overset{\sim}{\longrightarrow}\beta_*\OO_{\mathit{Spec}\mathscr{R}}$。

\textbf{17.1.C.（练习）}\quad 给定一个 $X$-态射 $\gamma$，
\begin{tikzcd}[row sep=2.2em, column sep=2.8em]
W \ar[rr, "\gamma"] \ar[dr, "\mu"'] && \mathit{Spec}\mathscr{R} \ar[dl, "\beta"] \\
& X, &
\end{tikzcd}
证明 \eqref{eq:relative-spec-univ-diagram} 中的那个态射 $\alpha$ 是合成
\[
\mathscr{R}\overset{\Phi}{\longrightarrow}\beta_*\OO_{\mathit{Spec}\mathscr{R}}\longrightarrow\beta_*\gamma_*\OO_W=\mu_*\OO_W.
\]

$\mathit{Spec}$ 构造给了我们一条理解\trAffine\trMorphism{}的重要途径。注意 $\mathit{Spec}\mathscr{R}\to X$ 是一个\trAffine\trMorphism。而「逆命题」也成立：

\textbf{17.1.D.（练习）}\quad 证明：若 $\mu\colon Z\to X$ 是一个\trAffine\trMorphism，则作为 $X$-\trScheme，我们有自然的\trIsomorphism{} $Z\overset{\sim}{\longleftarrow}\mathit{Spec}\mu_*\OO_Z$。

于是，借这个 $\mathit{Spec}$ 构造，我们可以复原任何\trAffine\trMorphism。更精确地说：一个\trMorphism{}是\trAffine{}的，当且仅当它具有上文所构造的 $\mu\colon\mathit{Spec}\mathscr{R}\to X$ 这种形式。

\textbf{17.1.E.（练习）}\quad 设 $\mu\colon\mathit{Spec}\mathscr{R}\to X$ 是这样一个\trMorphism。证明：$\mathit{Spec}\mathscr{R}$ 上诸\trQCSheaf{}的\trCategory{}等价于 $X$ 上带有 $\mathscr{R}$-\trModule{}结构的诸\trQCSheaf{}的\trCategory（即 $X$ 上拟凝聚的诸 $\mathscr{R}$-\trModule）。证明：$\mathit{Spec}\mathscr{R}$ 上有限型的诸\trQCSheaf{}之\trCategory{}等价于 $X$ 上带有「有限型」$\mathscr{R}$-\trModule{}结构的诸\trQCSheaf{}的\trCategory（即 $X$ 上有限型的拟凝聚 $\mathscr{R}$-\trModule）。

当 $X$ 十分简单而 $\mathit{Spec}\mathscr{R}$ 却颇为复杂时，这就很有用了。当 $X\cong\mathbb{P}^1$ 而 $\mathit{Spec}\mathscr{R}$ 是一条更复杂的\trCurve{}时，我们就会用到这一点；更一般地，在第 29 章证明\trSerreDuality{}时也会用到。

\textbf{17.1.F.（$\mathit{Spec}$ 的构造「与基变换交换」）}\quad 设 $\mu\colon Z\to X$ 是任意一个\trMorphism，而 $\mathscr{R}$ 是 $X$ 上一个\trQC{}\trSheafOfAlgebras。证明：存在自然的\trIsomorphism{} $Z\times_X\mathit{Spec}\mathscr{R}\overset{\sim}{\longleftarrow}\mathit{Spec}\mu^*\mathscr{R}$。（如果你琢磨一下，就会看到这推广了\mbox{命题~17.1.3} 的陈述，或许还有其证明。）

\textbf{17.1.5.}\ \emph{\trhDefinition{}。}\quad $\mathit{Spec}$ 构造的一个重要例子是\textbf{有限\trRank{}\trLocallyFreeSheaf{} $\mathscr{F}$ 的\trTotalSpace}（total space），我们把它定义为 $\mathit{Spec}(\operatorname{Sym}^\bullet\mathscr{F}^{\vee})$。

\textbf{17.1.G.（练习）}\quad 设 $\mathscr{F}$ 是\trRank{}为 $n$ 的一个\trLocallyFreeSheaf。证明：$\mathscr{F}$ 的\trTotalSpace{}是一个\trRank{}为 $n$ 的\emph{\trVectorBundle}，也就是说，给定任意一点 $p\in X$，都存在一个\trOpen\trNeighborhood{} $p\in\mathrm{U}\subset X$，使得
\[
\mathit{Spec}\bigl(\operatorname{Sym}^\bullet\mathscr{F}^{\vee}|_{\mathrm{U}}\bigr)\cong\mathbb{A}^n_{\mathrm{U}}.
\]
证明：$\mathscr{F}$ 同构于\trTotalSpace{} $\mathit{Spec}(\operatorname{Sym}^\bullet\mathscr{F}^{\vee})$ 的\trSection{}层。（可能的提示：使用\trTransitionFunction。）因此，\trTotalSpace{}也称作\textbf{与\trLocallyFreeSheaf{} $\mathscr{F}$ 相伴的\trVectorBundle}。（警示：一些作者，例如 [Stacks, tag 01M2]，把 $\mathit{Spec}(\operatorname{Sym}^\bullet\mathscr{F})$——即这个\trVectorBundle{}的\trDual——称作与 $\mathscr{F}$ 相伴的\trVectorBundle。）

特别地，若 $\mathscr{F}=\OO_X^{\oplus n}$，则 $\mathit{Spec}(\operatorname{Sym}^\bullet\mathscr{F}^{\vee})$ 就是 $\mathbb{A}^n_X$。

（旁注：你会注意到，构造 $\mathit{Spec}\operatorname{Sym}^\bullet$ 可以施用于任何\trCoherentSheaf{} $\mathscr{F}$（不取对偶，即 $\mathit{Spec}(\operatorname{Sym}^\bullet\mathscr{F})$）。这有时称作与 $\mathscr{F}$ 相伴的\textbf{\trAbelian\trCone}（abelian cone）。这个概念可能有用，但我们用不上它。）

\textbf{17.1.H.（$\mathbb{P}^n$ 上的\trTautologicalBundle{}是 $\OO(-1)$）}\quad 设 $k$ 是一个\trField。定义子集 $X\subset\mathbb{A}^{n+1}_k\times\mathbb{P}^n_k$ 为「$\mathbb{A}^{n+1}_k$ 中落在 $\mathbb{P}^n_k$ 的对应直线上的那些点」，使得映射 $\pi\colon X\to\mathbb{P}^n$ 在点 $\mathfrak{l}=[x_0,\cdots,x_n]$ 之上的\trFiber，恰是 $\mathbb{A}^{n+1}_k$ 中对应于 $\mathfrak{l}$ 的那条直线，也就是说，是 $(x_0,\dots,x_n)$ 的诸标量倍数。证明：$\pi\colon X\to\mathbb{P}^n_k$ 是（相应于\trInvertibleSheaf{} $\OO(-1)$ 的）那个\trLineBundle。（可能的提示：先在 $\mathbb{P}^n_k$ 的那些惯常的\trAffine{}开集上做，并弄清\trTransitionFunction。）因此，$\OO(-1)$ 常被称作 $\mathbb{P}^n_k$ 的\textbf{\trTautologicalBundle}（tautological bundle）（甚至在任意基上也是如此，而不只是域上）。（旁注：投射 $X\to\mathbb{A}^{n+1}_k$ 是 $\mathbb{A}^{n+1}_k$ 在「原点」处的\trBlowup；见\mbox{练习~10.3.F}。）

\textbf{17.1.6.\ 留待日后阅读：古怪的函子 $\pi^{!?}$，\trPushforward{}的一个偶现的右伴随。}\footnote{译注：上标「!?」整体是原书记号的一部分，并非排印笔误（页图放大核验）。}

由于接下来的论题会在两处不同的地方（\mbox{棘手练习~18.8.C} 与 \S29.3）被用到，而你或许会读、或许不会读它们，这段讨论除了放在这里，也再无其他自然的位置。但你现在不要去读它，而要等到需要它的时候再读，否则它将显得完全没有动机。

现在，当 $\pi$ 是一个\trFinite\trMorphism{}时，我们来为\trPushforward{} $\pi_*$ 定义一个\emph{右伴随} $\pi^{!?}$。（警示：记号 $\pi^{!?}$ 是非标准的，引入它只是出于我们将要给出的那些论证的需要。）这是令人惊讶的，因为我们通常把 $\pi_*$ 想成一个\emph{右伴随}（拉回 $\pi^*$ 的右伴随），而非左伴随。这可以粗略地类比于拉回那个出人意料的偶现\emph{左伴随}：\trExtensionByZero（\S23.4.7）。

\textbf{17.1.7.}\ \emph{\trhDefinition{}。}\quad 这个定义很简单，但很古怪。设 $\pi\colon X\to Y$ 是\trLocally\trNoetherian{}概形之间的一个\trFinite\trMorphism。若 $\mathscr{G}$ 是 $Y$ 上一个\trQCSheaf，则 $\mathscr{H}om_Y(\pi_*\OO_X,\mathscr{G})$ 也是 $Y$ 上一个\trQCSheaf（因为 $\pi_*\OO_X$ 是\trFinitelyPresented{}的，\mbox{练习~14.3.A(b)}）。它进一步带有 $(\pi_*\OO_X)$-\trModule{}的结构，这意味着它对应于 $X$ 上一个\trQCSheaf（由\mbox{练习~17.1.E}，该练习表明这个对应关系对 $\mathscr{G}$ 是函子性的）。令 $\pi^{!?}\mathscr{G}$ 为这个\trQCSheaf。那么 $\pi^{!?}$ 给出一个协变函子 $\mathit{QCoh}_Y\to\mathit{QCoh}_X$，并且也给出 $\mathit{Coh}_Y\to\mathit{Coh}_X$。

\textbf{17.1.I.（练习）}\quad 描述 $Y$ 上诸\trQCSheaf{}的一个自然\trIsomorphism
\[
\pi_*\pi^{!?}\mathscr{G}\overset{\sim}{\longrightarrow}\mathscr{H}om_Y(\pi_*\OO_X,\mathscr{G}),
\]
它对 $\mathscr{G}$ 是函子性的。等价地，描述函子的一个\trIsomorphism{} $\mathit{QCoh}_Y\to\mathit{QCoh}_Y$：
\[
\pi_*\pi^{!?}(\cdot)\overset{\sim}{\longrightarrow}\mathscr{H}om_Y(\pi_*\OO_X,\cdot).
\]

\textbf{17.1.8.}\quad 为了证明 $\pi^{!?}$ 是 $\pi_*$ 的右伴随，我们研究\trAffine{}情形，或者等价地，环论版本。设 $B\to A$ 是一个\trRingMorphism，$M$ 是一个 $A$-\trModule，$N$ 是一个 $B$-\trModule。注意 $\Hom_B(A,N)$ 自然带有 $A$-\trModule{}的结构。另外，$M$ 也自然带有 $B$-\trModule{}的结构；当我们想强调它作为 $B$-\trModule{}的结构时，有时便把它记作 $M_B$（见\mbox{练习~1.4.E}）。

考虑这个（十分奇怪的）映射
\begin{equation*}\tag{17.1.8.1}\label{eq:module-adjunction-map}
\Hom_A(M,\Hom_B(A,N))\longrightarrow\Hom_B(M_B,N)
\end{equation*}
其定义如下。给定 $m\in M$，以及 $\Hom_A(M,\Hom_B(A,N))$ 的一个元素 $\phi$，把 $m$ 送到 $\phi_m(1)$。

\textbf{17.1.J.（练习）}

（a）证明 \eqref{eq:module-adjunction-map} 是诸 $B$-\trModule{}之间的一个\trHomomorphism。

（b）证明 \eqref{eq:module-adjunction-map} 是一个双射。于是，$(M\mapsto M_B,\ N\mapsto\Hom_B(A,N))$ 构成一个\trAdjointPair{} $\mathit{Mod}_A\leftrightarrow\mathit{Mod}_B$。

\textbf{17.1.K.（练习）}\quad 设 $\pi\colon X\to Y$ 是\trLocally\trNoetherian{}概形之间的一个\trFinite\trMorphism。

（a）解释：前一练习中由 $N\mapsto\Hom_B(A,N)$ 给出的函子 $\mathit{Mod}_B\to\mathit{Mod}_A$ 如何就是\mbox{定义~17.1.7} 中 $\pi^{!?}$ 的\trAffineLocal{}描述。

（b）证明 $(\pi_*,\pi^{!?})$ 构成 $\mathit{QCoh}_X$ 与 $\mathit{QCoh}_Y$ 之间的一个\trAdjointPair。

（c）证明 $(\pi_*,\pi^{!?})$ 构成 $\mathit{Coh}_X$ 与 $\mathit{Coh}_Y$ 之间的一个\trAdjointPair。

（d）若 $\mathscr{F}\in\mathit{QCoh}_X$ 且 $\mathscr{G}\in\mathit{QCoh}_Y$，描述一个自然\trIsomorphism
\begin{equation*}\tag{17.1.8.2}\label{eq:sheaf-adjunction-isomorphism}
\pi_*\mathscr{H}om_X(\mathscr{F},\pi^{!?}\mathscr{G})\overset{\sim}{\longrightarrow}\mathscr{H}om_Y(\pi_*\mathscr{F},\mathscr{G})
\end{equation*}
（对 $\mathscr{F}$ 与 $\mathscr{G}$ 都是函子性的），它在\trAffineLocal{}上恰是\mbox{练习~17.1.J} 中所描述的那个\trIsomorphism{} \eqref{eq:module-adjunction-map}。


\section{（拟凝聚）分次代数层的相对 $\Proj$}

与 $\Spec$ 的相对版本 $\mathit{Spec}$ 相平行，我们对概形 $X$ 上一个拟凝聚的\trSheafOfGradedAlgebras（在满足某些假设之下）定义 $\Proj$ 的一个相对版本，记作 $\mathit{Proj}$（称作「相对 $\Proj$」或「层 $\Proj$」）。基 $X$ 为仿射的情形，我们已经在 \S4.5.7 中处理过了，那里用的是环 $A$ 上的常规 $\Proj$ 构造。优雅的推进办法，本是先陈述恰当的\trUniversalProperty，然后聪明地利用它把各个仿射片上的构造粘合起来，正如我们在\trFiberProduct、\trNormalization{}与 $\mathit{Spec}$ 的构造中所做的那样。但由于\trGradedRing{}和\trGradedModule{}会把一切都搅得令人糊涂，我们并不这样做，而是心虚地采取一条更为平淡无奇的路线。\footnote{译注：pedestrian 本义为「步行的」，引申为「平淡乏味、缺乏想象力」，作者在此自嘲这条路线缺乏格调。}（这个\trUniversalProperty{}是可以做成的，见 [Stacks, tag 01O0]。但我建议你莫这么做——一个\trUniversalProperty{}应当让你的日子更好过，而不是更难过。）

\textbf{17.2.A.（练习：$\Proj$ 与仿射基变换交换）}\quad 设 $A\to B$ 是环映射\footnote{原书此处作「is map of rings」，当为「is a map of rings」之脱漏，译文径顺化。——译者注}，而 $S_\bullet$ 是 $A$ 上一个 $\mathbb{Z}^{\ge0}$-\trGradedRing。

（a）给出一个典范的\trIsomorphism
\begin{equation*}\tag{17.2.0.1}\label{eq:relproj-affine-base-change}
\alpha\colon\Proj_B(S_\bullet\otimes_A B)\overset{\sim}{\longrightarrow}(\Proj_A S_\bullet)\times_{\Spec A}\Spec B.
\end{equation*}
（$\Proj$ 上的\textbf{这些下标}是一种新记号。它们是多余的，只不过是提醒我们，比如说，此刻我们是把 $\Proj_A S_\bullet$ 当作一个 $A$-\trScheme{}来想的。）

（b）（容易）设 $X$ 是一个射影 $A$-\trScheme（\S4.5.10）。证明：$X\times_{\Spec A}\Spec B$ 是一个射影 $B$-\trScheme。

（c）设 $S_\bullet$ 在次数 $1$ 处生成，于是 $\OO_{\Proj_A S_\bullet}(1)$ 是一个\trInvertibleSheaf（\mbox{练习~15.1.G}）。显然，$S_\bullet\otimes_A B$ 作为一个 $B$-\trAlgebra{}也在次数 $1$ 处生成。描述一个\trIsomorphism
\[
\OO_{\Proj_B(S_\bullet\otimes_A B)}(1)\overset{\sim}{\longleftarrow}\alpha^*\gamma^*\OO_{\Proj_A S_\bullet}(1),
\]
其中 $\gamma$ 是如下\trCartesianDiagram{}中上方的那条\trMorphism：
\[
\begin{tikzcd}
(\Proj_A S_\bullet)\times_{\Spec A}\Spec B \ar[r, "\gamma"] \ar[d] & \Proj_A S_\bullet \ar[d] \\
\Spec B \ar[r] & \Spec A.
\end{tikzcd}
\]
可能的提示：\trTransitionFunction。

现在我们给出一个在 $X$ 上构造概形的通用办法（取自 [Stacks, tag 01LH]），它适用于这样的场合：我们知道这些概形在任何\trAffine{}开集上应当是什么，也知道它们在从一个\trAffine{}开集开嵌入另一个\trAffine{}开集之时如何表现。

\textbf{17.2.B.（练习）}\quad 设给定一个\trScheme{} $X$，以及如下数据：

（i）对每个\trAffine{}开子集 $U\subset X$，给定某个\trMorphism{} $\pi_U\colon Z_U\to U$（「$U$ 上的一个概形」）。

（ii）对仿射开子集的每个（开）包含 $V\subset U\subset X$，给定一个\trOpenEmbedding{} $\rho^U_V\colon Z_V\hookrightarrow Z_U$。

假设这份数据满足：

（a）对每个 $V\subset U\subset X$，$\rho^U_V$ 诱导出 $V$ 上概形的一个\trIsomorphism{} $Z_V\overset{\sim}{\longrightarrow}\pi_U^{-1}(V)$；并且

（b）只要 $W\subset V\subset U\subset X$ 是三个依次相嵌的\trAffine{}开子集，就有 $\rho^U_W=\rho^U_V\circ\rho^V_W$。

证明：存在一个 $X$-\trScheme{} $\pi\colon Z\to X$，并且在每个\trAffine{}开集 $U$ 上存在\trIsomorphism{} $i_U\colon\pi^{-1}(U)\overset{\sim}{\longrightarrow}Z_U$，使得对依次相嵌的\trAffine{}开集 $V\subset U$，$\rho^U_V$ 与下述复合一致：
\[
\begin{tikzcd}
Z_V \ar[r, "i_V^{-1}"] & \pi^{-1}(V) \ar[r, hook] & \pi^{-1}(U) \ar[r, "i_U"] & Z_U.
\end{tikzcd}
\]
提示（参\mbox{练习~4.4.A}）：先把 $Z$ 作为\trSet{}构造出来，再作为\trTopologicalSpace{}构造出来，最后作为\trScheme{}构造出来。（你的构造将不依赖于种种选择。你的解法在更一般的情形也会奏效，比如说，把概形的\trCategory{}换成赋环空间，把\trAffine{}开子集换成拓扑的任意一组\trBase。）

\textbf{17.2.C.（重要练习与定义：相对 $\mathit{Proj}$）}\quad 设 $\mathscr{S}_\bullet=\oplus_{n\ge0}\mathscr{S}_n$ 是概形 $X$ 上一个拟凝聚的 $\mathbb{Z}^{\ge0}$-\trSheafOfGradedAlgebras。在每个\trAffine{}开子集 $\Spec A\cong U\subset X$ 上，我们有一个 $U$-\trScheme{} $\Proj_A\mathscr{S}_\bullet(U)\to U$。证明：它们可以粘合起来，形成一个 $X$-\trScheme，我们称之为 $\mathit{Proj}_X\mathscr{S}_\bullet$；我们有一条「\trStructureMorphism」$\beta\colon\mathit{Proj}_X\mathscr{S}_\bullet\to X$。（\trStructureMorphism{} $\beta$ 是定义的一部分。）提示／警告：倘若\mbox{练习~17.2.B} 只要求考虑\emph{特异的}仿射包含，这个问题会更容易；但那又会使\mbox{练习~17.2.B} 明显变难。所以，请改用\mbox{练习~17.2.A}。

由\mbox{练习~17.2.B} 的构造，任何\trAffine{}开集之上的原像都可以用原本的 $\Proj$ 构造算出。（你或许乐于回头用这个想法，给出 $X^{\mathrm{red}}$、$X$ 的\trNormalization，以及一个拟凝聚 $\OO$-\trAlgebra{}层的 $\mathit{Spec}$ 的构造。但放弃\trUniversalProperty{}，是要付出道德代价的。）

\textbf{17.2.1.}\ \textbf{关于 $\mathscr{S}_\bullet$ 的常设（合理）假设：「在次数 $1$ 处有限生成」。}\quad $\Proj$ 构造最有用的场合，是把它施加于一个满足某些合理假设的 $A$-\trAlgebra{} $S_\bullet$（\S4.5.6）；特别是当 $S_\bullet$ 是一个\trFinitelyGenerated{}的 $\mathbb{Z}^{\ge0}$-分次 $A$-\trAlgebra{}时，而若它还在次数 $1$ 处生成，则最为理想。因此，在本书余下的部分里，一旦我们对拟凝聚\trSheafOfAlgebras{}把这些假设说清楚，我们就会对 $\mathscr{S}_\bullet$ 一律强制执行它们。（如果你日后需要放宽这些假设——比如说，保留有限生成假设而去掉「在次数 $1$ 处生成」假设——也不会太困难。）确切地说，\textbf{我们现在总是要求：（i）$\mathscr{S}_0=\OO_X$；（ii）$\mathscr{S}_\bullet$ 「在次数 $1$ 处生成」；（iii）$\mathscr{S}_1$ 是\trFiniteType{}的。}把条件（ii）说精确的最干净方式，是要求自然映射
\[
\Sym^\bullet_{\OO_X}\mathscr{S}_1\longrightarrow\mathscr{S}_\bullet
\]
为\trSurjectiveMap。由于 $\Sym^\bullet$ 构造可以仿射局部地计算（\S14.2.6），我们可以在任何\trAffine\trCover{}上检验「在次数 $1$ 处生成」。这些假设的动机，来自下面这个构造。

\textbf{17.2.D.（重要练习：$\mathit{Proj}_X\mathscr{S}_\bullet$ 上的 $\OO(1)$）}\quad 若 $\mathscr{S}_\bullet$ 在次数 $1$ 处有限生成（\mbox{假设~17.2.1}），构造 $\mathit{Proj}_X\mathscr{S}_\bullet$ 上的一个\trInvertibleSheaf{} $\OO_{\mathit{Proj}_X\mathscr{S}_\bullet}(1)$，使之「在每个\trAffine{}开子集 $\Spec A\subset X$ 上，限制为 $\OO_{\Proj_A\mathscr{S}_\bullet(\Spec A)}(1)$」。

\textbf{17.2.E.（练习：「$\mathit{Proj}$ 与基变换交换」）}\quad 设 $\mathscr{S}_\bullet$ 是 $X$ 上一个拟凝聚的 $\mathbb{Z}^{\ge0}$-\trSheafOfGradedAlgebras。令 $\rho\colon Z\to X$ 为任一\trMorphism。给出一个自然的\trIsomorphism
\[
(\mathit{Proj}\,\rho^*\mathscr{S}_\bullet\,,\ \OO_{\mathit{Proj}\,\rho^*\mathscr{S}_\bullet}(1))\overset{\sim}{\longleftarrow}(Z\times_X\mathit{Proj}\,\mathscr{S}_\bullet\,,\ \psi^*\OO_{\mathit{Proj}\,\mathscr{S}_\bullet}(1)),
\]
其中 $\psi$ 是基变换图中「上方」的那条\trMorphism：
\[
\begin{tikzcd}
Z\times_X\mathit{Proj}\,\mathscr{S}_\bullet \ar[r, "\psi"] \ar[d] & \mathit{Proj}\,\mathscr{S}_\bullet \ar[d, "\beta"] \\
Z \ar[r, "\rho"] & X.
\end{tikzcd}
\]

\textbf{17.2.2.}\ \emph{\trhDefinition：$\pi$-极丰富。}\quad 设 $\pi\colon X\to Y$ 是\trProper{}的。若 $\mathscr{L}$ 是 $X$ 上的一个\trInvertibleSheaf，则我们说 $\mathscr{L}$ 是\textbf{相对于 $\pi$ 极丰富}的、\textbf{相对极丰富}的，或者（说得很别扭）\textbf{$\pi$-极丰富}的，若我们能写成 $X=\mathit{Proj}_Y\mathscr{S}_\bullet$ 且 $\mathscr{L}\cong\OO(1)$，其中 $\mathscr{S}_\bullet$ 是 $Y$ 上一个满足\mbox{假设~17.2.1}（「在次数 $1$ 处有限生成」）的拟凝聚\trSheafOfAlgebras。（\trVeryAmpleness{}的概念可以推广到更一般的情形；例如见 [Stacks, tag 01VM]。但这只对趣味冷僻的人才有意思。）

\textbf{17.2.F.（练习）}\quad 设 $\mathscr{S}_\bullet$ 在次数 $1$ 处有限生成（\mbox{假设~17.2.1}）。描述分次\trQCSheaf{}之间的一个映射 $\phi\colon\mathscr{S}_\bullet\to\oplus_{n=0}^{\infty}\beta^*\OO(n)$（$\beta$ 是\trStructureMorphism；见\mbox{练习~17.2.C}）。提示：\mbox{练习~15.7.C}。

\textbf{17.2.G.（练习）}\quad 设 $\mathscr{L}$ 是 $X$ 上的一个\trInvertibleSheaf，而 $\mathscr{S}_\bullet$ 是 $X$ 上一个在次数 $1$ 处生成的拟凝聚\trSheafOfGradedAlgebras（\mbox{假设~17.2.1}）。定义 $\mathscr{S}_\bullet'=\oplus_{n=0}^{\infty}(\mathscr{S}_n\otimes\mathscr{L}^{\otimes n})$。于是 $\mathscr{S}_\bullet'$ 从 $\mathscr{S}_\bullet$ 那里继承了一个自然的代数结构；请描述它。给出「带线丛的 $X$-概形」之间的一个自然\trIsomorphism
\[
(\mathit{Proj}\,\mathscr{S}_\bullet'\,,\ \OO_{\mathit{Proj}\,\mathscr{S}_\bullet'}(1))\overset{\sim}{\longleftrightarrow}(\mathit{Proj}\,\mathscr{S}_\bullet\,,\ \OO_{\mathit{Proj}\,\mathscr{S}_\bullet}(1)\otimes\beta^*\mathscr{L}),
\]
其中 $\beta\colon\mathit{Proj}\,\mathscr{S}_\bullet\to X$ 是\trStructureMorphism。换句话说，撇开形式地说：$\mathit{Proj}$ 还是原来那一个，只是 $\OO(1)$ 被 $\mathscr{L}$ 扭转了一下。

\textbf{17.2.3.}\ \textbf{\trVectorBundle{}的\trProjectivization。}

\textbf{17.2.4.}\ \emph{\trhDefinition。}\quad 若 $\mathscr{F}$ 是 $X$ 上一个\textbf{有限秩的\trLocallyFreeSheaf}，则 $\mathit{Proj}(\Sym^\bullet\mathscr{F}^\vee)$ 称为它的\textbf{\trProjectivization}（projectivization），并记作 $\mathbb{P}\mathscr{F}$。（取对偶的理由，与\mbox{定义~17.1.5} 中取 $\mathit{Spec}(\Sym^\bullet\mathscr{F}^\vee)$ 的理由相同。）你可以自行检验，这一构造在\trBaseChange{}下表现良好。定义 $\mathbb{P}^n_X:=\mathbb{P}(\OO_X^{\oplus(n+1)})$。（于是 $\mathbb{P}^n_{\Spec A}$ 与我们早先对 $\mathbb{P}^n_A$ 的定义相符——参\mbox{练习~4.5.Q}——而 $\mathbb{P}^n_X$ 与我们早先的用法相符；例如见\mbox{定理~11.4.5} 的证明。）若 $\mathscr{F}$ 是秩为 $n+1$ 的\trLocallyFreeSheaf，则 $\mathbb{P}\mathscr{F}$ 是 $X$ 上的一束\textbf{射影丛}或 $\mathbb{P}^n$-\textbf{丛}。由\mbox{练习~17.2.G}，若 $\mathscr{F}$ 是 $X$ 上一个有限秩的\trLocallyFreeSheaf，则存在典范的\trIsomorphism{} $\mathbb{P}\mathscr{F}\overset{\sim}{\longleftarrow}\mathbb{P}(\mathscr{L}\otimes\mathscr{F})$。

更一般地，若 $\mathscr{F}$ 是 $X$ 上一个\trFiniteType{}的\trQCSheaf，那么也可以类似地定义它的\trProjectivization{} $\mathit{Proj}(\Sym^\bullet\mathscr{F}^\vee)$。不过要当心。例如，若 $\mathscr{G}$ 是某个整概形上的\trTorsion\trSheaf，则 $\mathscr{F}^\vee=0$\footnote{原书此句先设 $\mathscr{G}$ 为挠层，随后却作「$\mathscr{F}^\vee=0$」与「$\mathbb{P}\mathscr{F}=\varnothing$」，疑当为 $\mathscr{G}$，兹照录原书。——译者注}，于是按这个定义便有 $\mathbb{P}\mathscr{F}=\varnothing$。所以这并不是一个多好的概念。

\textbf{17.2.5.}\ \emph{记号的特大警示。}\quad 关于 $\mathbb{P}\mathscr{F}$ 究竟应当定义为 $\Sym^\bullet\mathscr{F}$ 还是 $\Sym^\bullet\mathscr{F}^\vee$，各家势同水火，正如所谓「向量空间的\trProjectivization{}到底参数化一维\trSubspace{}还是一维商」的说法（参\mbox{练习~17.2.I}）一样。因此，最稳妥的办法是避开记号 $\mathbb{P}\mathscr{F}$，或者至少开宗明义地声明你遵循的是哪一种约定。\emph{在本书中，我们将尽量避免使用记号 $\mathbb{P}\mathscr{F}$，并且只在 $\mathscr{F}$ 是有限秩局部自由层（即「一个\trVectorBundle」）时才使用它，而它的含义将永远指这里所定义的 $\mathit{Proj}(\Sym^\bullet\mathscr{F}^\vee)$。}

\textbf{17.2.6.}\ \emph{例：\trRuling{}曲面。}\quad 若 $C$ 是一条\trRegular\trCurve，而 $\mathscr{F}$ 是秩 $2$ 的\trLocallyFreeSheaf，则 $\mathbb{P}\mathscr{F}$ 称为 $C$ 上的一张\textbf{\trRuling{}曲面}（ruled surface）。若 $C$ 还\trIsomorphic{}于 $\mathbb{P}^1$，则 $\mathbb{P}\mathscr{F}$ 称为一张\textbf{\trHirzebruchSurface}（Hirzebruch surface）。$\mathbb{P}^1$ 上的所有\trVectorBundle{}都分裂为\trLineBundle{}的\trDirectSum（证明见 \S18.5.5），故每张\trHirzebruchSurface{}都形如 $\mathbb{P}(\OO(n_1)\oplus\OO(n_2))$。由\mbox{练习~17.2.G}，这只依赖于 $n_2-n_1$。\trHirzebruchSurface{} $\mathbb{P}(\OO\oplus\OO(n))$（$n\ge0$）常记作 $\mathbb{F}_n$。我们将在 \S20.2.10 中用更长的篇幅讨论\trHirzebruchSurface{}；在\mbox{练习~20.2.Q} 中我们将会看到，诸 $\mathbb{F}_n$ 两两互异。

\textbf{17.2.H.（练习）}\quad 若 $\mathscr{S}_\bullet$ 在次数 $1$ 处有限生成（\mbox{假设~17.2.1}），描述一个典范的\trClosedEmbedding
\[
\begin{tikzcd}
\mathit{Proj}\,\mathscr{S}_\bullet \ar[rr, hook, "i"] \ar[dr, "\beta"'] && \mathit{Proj}_X\Sym^\bullet\mathscr{S}_1 \ar[dl] \\
& X &
\end{tikzcd}
\]
以及一个\trIsomorphism{} $\OO_{\mathit{Proj}\,\mathscr{S}_\bullet}(1)\overset{\sim}{\longleftarrow}i^*\OO_{\mathit{Proj}_X\Sym^\bullet\mathscr{S}_1}(1)$，它由满射
\[
\Sym^\bullet\mathscr{S}_1\twoheadrightarrow\mathscr{S}_\bullet.
\]
诱导而来。特别地，若 $\mathscr{S}_1$ 是\trLocallyFree{}的，则 $\mathit{Proj}\,\Sym^\bullet\mathscr{S}_1=\mathbb{P}\mathscr{S}_1^\vee$，于是我们便把 $\mathit{Proj}\,\mathscr{S}_\bullet$ 嵌入了 $X$ 上的一束射影丛。

\textbf{17.2.I.（练习）}\quad 设 $\mathscr{F}$ 是 $X$ 上秩为 $n+1$ 的\trLocallyFreeSheaf。在「$\mathbb{P}\mathscr{F}\to X$ 的\trSection{} $s\colon X\to\mathbb{P}\mathscr{F}$」全体所成之集，与「$\mathscr{F}$ 到 $X$ 上诸\trInvertibleSheaf{}的满同态 $\mathscr{F}\to\mathscr{L}\to 0$」全体所成之集之间，给出一个双射。$\mathbb{P}\mathscr{F}$ 的这一函子式描述，在某种意义上推广了 \S15.2.3 中对\trProjectiveSpace{}的函子式描述。

\textbf{17.2.7.}\ \emph{\trhRemark{}（\trProjective\trCone{}与\trAffine\trCone{}的相对版本）。}\quad 从 $\mathit{Spec}\,\mathscr{S}_\bullet$ 去掉零截面之后，有一个到 $\mathit{Proj}\,\mathscr{S}_\bullet$ 的自然\trMorphism（参\mbox{练习~9.3.N}）。正如 $\Proj S_\bullet[T]$ 含有一个可与 $\Proj S_\bullet$ 相认同的\trClosedSubscheme、其补集可与 $\Spec S_\bullet$ 相认同（\mbox{练习~9.3.O}），$\mathit{Proj}\,\mathscr{S}_\bullet[T]$ 也含有一个可与 $\mathit{Proj}\,\mathscr{S}_\bullet$ 相认同的\trClosedSubscheme，其补集可与 $\mathit{Spec}\,\mathscr{S}_\bullet$ 相认同。欢迎你把这件事想清楚。


\section{射影态射}

在 \S17.1 中，我们重新诠释了\trAffine\trMorphism：$X\to Y$ 是一个\trAffine\trMorphism，若对 $Y$ 上某个\trQC{}\trSheafOfAlgebras{} $\mathscr{B}$，存在 $Y$-概形之间的\trIsomorphism{} $X\overset{\sim}{\longleftarrow}\mathit{Spec}\,\mathscr{B}$。我们将以类似的方式来\emph{定义}\trProjective\trMorphism{}的观念。

你或许会觉得，射影性既然是如此经典的一个概念，总该存在某种表现合理的显然定义才是。但情况并非如此：射影的可能的变式定义有很多（见 [Stacks, tag 01W8]），而且它们全都有缺陷，包括我们在此给出的那个公认定义在内（见 \S17.3.4 中的种种警示）。

\textbf{17.3.1.}\ \emph{\trhDefinition{}。}\quad 一个\trMorphism{} $\pi\colon X\to Y$ 称为\textbf{\trProjective{}的}（projective），若存在一个\trIsomorphism
\[
\begin{tikzcd}[row sep=2.4em, column sep=4.6em]
X \ar[rr, "\sim"] \ar[dr, "\pi"'] && \mathit{Proj}\,\mathscr{S}_\bullet \ar[dl] \\
& Y &
\end{tikzcd}
\]
其中 $\mathscr{S}_\bullet$ 是 $Y$ 上一个满足「在次数 $1$ 处有限生成」（\mbox{假设~17.2.1}）的\trQC{}\trSheafOfAlgebras。我们称 $X$ 为一个\textbf{\trProjective{} $Y$-\trScheme}（projective $Y$-scheme），或者说 $X$ \textbf{在 $Y$ 上是\trProjective{}的}（projective over $Y$）。利用\mbox{练习~7.4.D}，这推广了射影 $A$-\trScheme{}的观念。

\textbf{17.3.A.（练习）}

（a）\emph{（射影态射的一个有用刻画）}设 $\pi\colon X\to Y$ 是一个\trMorphism。证明：$\pi$ 是射影的，当且仅当存在 $Y$ 上一个\trFiniteType{}的\trQCSheaf{} $\mathscr{S}_1$，以及一个\trClosedEmbedding{} $\iota\colon X\hookrightarrow\mathit{Proj}_Y\Sym^\bullet\mathscr{S}_1$（在 $Y$ 上，即与到 $Y$ 的诸\trMap{}相交换）。提示：\mbox{练习~17.2.H}。

（b）\emph{（射影态射的一个有用刻画，带线丛）}设 $\mathscr{L}$ 是 $X$ 上的一个\trInvertibleSheaf，而 $\pi\colon X\to Y$ 是一个\trMorphism。证明：$\pi$ 是射影的、且 $\OO(1)\cong\mathscr{L}$，当且仅当存在 $Y$ 上一个\trFiniteType{}的\trQCSheaf{} $\mathscr{S}_1$、一个\trClosedEmbedding{} $\iota\colon X\hookrightarrow\mathit{Proj}_Y\Sym^\bullet\mathscr{S}_1$（在 $Y$ 上，即与到 $Y$ 的诸\trMap{}相交换），以及一个\trIsomorphism{} $\iota^*\OO_{\mathit{Proj}_Y\Sym^\bullet\mathscr{S}_1}(1)\overset{\sim}{\longleftarrow}\mathscr{L}$。

（c）再进一步设 $Y$ 容许一个 \S16.2.5 意义下的\trAmple\trLineBundle——每当 $Y$ 是射影的、仿射的，或更一般地是拟射影的时，情形便是如此。证明：$\pi$ 是射影的，当且仅当对某个 $n$，存在一个\trClosedEmbedding{} $X\to\mathbb{P}^n_Y$（在 $Y$ 上）。这正是 [Ha1, p.~103] 中所给出的\trProjective\trMorphism{}定义。（如果你想避开带星的 \S16.2.5，那么可以假定 $Y$ 在 $\Spec A$ 上是射影的，并使用 \S16.2.1 中丰富性的定义。这样你便处理了 $Y$ 为射影的这一重要情形，但会错过其他可能有趣的情形，比如 $Y$ 为仿射或拟射影（但不\trProper）的情形。）提示：较难的方向是正向蕴含。用（a）中那个\trFiniteType{}的\trQCSheaf{} $\mathscr{S}_1$。把 $\mathscr{S}_1$ 与 $\mathscr{M}$ 的一个足够高的幂作张量，使之有限整体生成（\mbox{定理~16.2.6}，或在\trProper{}设定下的\mbox{定理~16.2.2}），从而得到一个满射
\[
\OO_Y^{\oplus(n+1)}\twoheadrightarrow\mathscr{S}_1\otimes\mathscr{M}^{\otimes\mathrm{N}}.
\]
然后用\mbox{练习~17.2.G}。

\textbf{17.3.2.}\ \textbf{射影态射的性质。}

一个态射之为射影的这一性质，显然在\trBaseChange{}下保持，因为 $\mathit{Proj}$ 构造对\trBaseChange{}表现良好（\mbox{练习~17.2.E}）。此外，\trProjective\trMorphism{}是\trProper{}的：\trProperness{}在靶上是局部的（\mbox{定理~11.4.4}(b)），而我们此前已见到射影 $A$-\trScheme{}在 $A$ 上是\trProper{}的（\mbox{定理~11.4.5}）。特别地（由\trProperness{}的定义），\trProjective\trMorphism{}是\trSeparated{}的、\trFiniteType{}的，并且是\trUniversallyClosed{}的。

\textbf{17.3.B.（重要练习：有限态射是射影的；对照\mbox{练习~8.3.1}）}\quad 按如下思路证明有限态射是射影的。设 $Z\to X$ 是\trFinite{}的，于是 $Z\cong\mathit{Spec}\,\mathscr{B}$，其中 $\mathscr{B}$ 是 $X$ 上一个\trFiniteType{}的\trQC{}\trSheafOfAlgebras。描述一个\trSheafOfGradedAlgebras{} $\mathscr{S}_\bullet$，使得 $\mathscr{S}_0\cong\OO_X$，而对 $n>0$ 有 $\mathscr{S}_n\cong\mathscr{B}$。描述一个 $X$-\trIsomorphism{} $Z\overset{\sim}{\longleftarrow}\mathit{Proj}\,\mathscr{S}_\bullet$。

特别地，\trClosedEmbedding{}是射影的。对于态射，我们有如下的蕴含序列
\[
\text{\trClosedEmbedding}\Longrightarrow\text{\trFinite}\Longrightarrow\text{\trProjective}\Longrightarrow\text{\trProper}.
\]

\textbf{17.3.3.}\quad 我们已经证明了有限态射是射影的（\mbox{练习~17.3.B}），并且具有有限的\trFiber（\mbox{练习~8.3.J}）。我们将在\mbox{定理~18.1.6} 中证明其逆，并在其后立刻陈述向\trProper\trMorphism{}的推广（尽管我们要到\mbox{定理~28.5.2} 才会证明它）。

\textbf{17.3.4.}\ \emph{关于射影态射的种种警示。}

首先，注意 $\OO(1)$——$\mathit{Proj}$ 这个概念中的一个重要组成部分——在定义中并没有被提到。（我倒宁愿它是定义的一部分，但这并不是公认的做法。）

其次，[Ha1, p.~103] 对\trProjective\trMorphism{}给出了一个不同的定义；我们遵循的是 \trGrothendieck{}的那个更一般的定义。尽管这些定义在多数场合下是同一回事（\mbox{练习~17.3.A}(c)），你在阅读文献时仍应保持警惕。例如，有限态射在 [Ha1] 的意义下并不总是射影的，而按这里的定义它们\emph{却}是射影的（\mbox{练习~17.3.B}）。

第三，射影态射在 \S8.1 的意义下是\emph{不}\trReasonable{}的！尽管\trProjective\trMorphism{}在\trBaseChange{}下保持（\S17.3.2），「是射影的」这一性质在靶上却\emph{不}是局部的（例子见 \S24.7.7）。（若再添加\trInvertibleSheaf{} $\OO(1)$ 这份额外数据，并且在诺特的情形下，它便是局部的了；见 \S17.3.8。）此外，我们将会证明，仅当最终的靶是仿射或\trNoetherian{}的时，射影态射才在复合下保持（\mbox{练习~17.3.I}）。

但是，合理性（reasonableness）的诸多推论中有许多依然成立，正如接下来几道练习所表明的那样。我们首先要证明的是：射影态射这一类态射，尽管本性不合理，却仍然满足\trCancellationTheorem~11.1.1。

\textbf{17.3.C.（练习：射影态射的消去定理）}

（a）设 $\pi\colon X\to Y$ 是一个\trClosedEmbedding，而 $\rho\colon Y\to Z$ 是射影的。证明：$\rho\circ\pi$ 是射影的。

（b）证明：一个\trMorphism{} $\rho\colon Y\to Z$ 具有射影对角，当且仅当 $\rho$ 是\trSeparated{}的。

（c）证明射影态射的消去定理：设 $\pi\colon X\to Y$ 是一个\trMorphism，$\rho\colon Y\to Z$ 是一个\trSeparated\trMorphism（亦即，由（b），它具有射影对角），而 $\rho\circ\pi$ 是射影的。证明：$\pi$ 是射影的。（可能的提示：使用\mbox{练习~11.1.A} 的提示中的那个图。）

\textbf{17.3.D.（练习）}\quad 证明：从一个射影 $k$-\trScheme{}到一个\trSeparated{} $k$-\trScheme{}的（$\Spec k$ 上的）\trMorphism{}总是射影的。（提示：前一练习。）

\textbf{17.3.E.（练习）}\quad 设 $\pi\colon X\to Y$ 与 $\pi'\colon X'\to Y'$ 是 $S$-\trScheme{}之间的两个射影态射。证明：$\pi\times_S\pi'\colon X\times_S X'\to Y\times_S Y'$ 也是射影的。提示：鉴于射影态射的不合理本性，你不能求助于\mbox{练习~8.1.A}；请改用\mbox{练习~10.6.D}。

\textbf{17.3.5.}\ \hard\ \textbf{相对情形下的\trGlobalGeneration{}与（极）\trAmpleness。}

我们把 \S16.1 的讨论推广到相对情形，为的是给自己配上相对无基点性（relatively base-point-freeness）的语言。我们不会用到这一讨论，所以在初读时你应当跳过它。但这些观念在研究文献中会反复出现。

设 $\pi\colon X\to Y$ 是一个\trQuasicompact\trQuasiSeparated\trMorphism。若 $\mathscr{F}$ 是 $X$ 上一个\trQCSheaf，我们就说 $\mathscr{F}$ 是\textbf{相对\trGloballyGenerated}（relatively globally generated）或\textbf{相对于 $\pi$ \trGloballyGenerated}（globally generated with respect to $\pi$），若拟凝聚层之间的自然映射 $\pi^*\pi_*\mathscr{F}\to\mathscr{F}$ 是满射。要保证 $\pi_*\mathscr{F}$ 是一个\trQCSheaf，拟紧性与拟分离性是必需的\footnote{原书此处作「are needed ensure」，当为「are needed to ensure」之脱漏，译文径顺化。——译者注}（\mbox{练习~14.4.C}）。但这些假设并不算很强的限制。整体生成只在拟紧的设定中才最有用，而多数人也不会为拟分离假设所烦扰。（不重要的旁注：这些假设其实可以大幅放宽。设 $\pi\colon X\to Y$ 是\emph{\trLocallyRingedSpace}之间的一个\trMorphism——不必是概形——而不带任何其他假设，且 $\mathscr{F}$ 是 $X$ 上一个\trQCSheaf，那么若 $\OO_X$-模之间的自然映射 $\pi^*\pi_*\mathscr{F}\to\mathscr{F}$ 是满射，我们就说 $\mathscr{F}$ 是相对\trGloballyGenerated{}或相对于 $\pi$ \trGloballyGenerated。）

得益于我们的这些假设，自然映射 $\pi^*\pi_*\mathscr{F}\to\mathscr{F}$ 是拟凝聚层之间的一个\trMorphism，于是「相对整体生成」这一条件在 $Y$ 上是仿射局部的。

现在设 $\mathscr{L}$ 是 $X$ 上一个\trLocallyFreeSheaf，而 $\pi\colon X\to Y$ 是一个\trMorphism。我们说 $\mathscr{L}$ 是\textbf{相对\trBasePointFree}的（relatively base-point-free）或\textbf{相对于 $\pi$ \trBasePointFree}的（base-point-free with respect to $\pi$），若它是相对\trGloballyGenerated。

\textbf{17.3.F.（练习）}\quad 设 $\mathscr{L}$ 是 $X$ 上一个相对\trBasePointFree{}的\trLineBundle，$\pi\colon X\to Y$ 是一个\trQuasicompact\trSeparated\trMorphism，而 $\pi_*\mathscr{L}$ 在 $Y$ 上是\trFiniteType{}的。（我们日后将在 \trGrothendieck{}凝聚定理~18.9.1 中证明：若 $\pi$ 是\trProper{}的且 $Y$ 是\trNoetherian{}的，则后一陈述为真。当 $\pi$ 是射影的时，这要容易得多；见\mbox{定理~18.7.1}。我们本可以现在就下功夫证明它，但那不值得费事。）描述一个典范\trMorphism{} $\psi\colon X\to\mathit{Proj}_Y\Sym^\bullet(\pi_*\mathscr{L})$。（可能的提示：这推广了「无基点的线丛给出到射影空间的映射」这一事实，所以请把那个论证加以推广；见 \S15.2。）

我们说 $\mathscr{L}$ 是\textbf{相对\trAmple}的（relatively ample）、\textbf{$\pi$-\trAmple}的（$\pi$-ample），或者\textbf{相对于 $\pi$ \trAmple}的（relatively ample with respect to $\pi$），若对 $Y$ 的每个\trAffine{}开子集 $\Spec B$，$\mathscr{L}|_{\pi^{-1}(\Spec B)}$ 在 $\pi^{-1}(\Spec B)$ 上于 $B$ 上是\trAmple{}的；或者，等价地（由 \S16.2.5），$\mathscr{L}|_{\pi^{-1}(\Spec B)}$ 在 $\pi^{-1}(\Spec B)$ 上是（绝对）\trAmple{}的。由 \S16.2.5 的讨论，若 $\mathscr{L}$ 是\trAmple{}的\footnote{原书此处作「if $\mathscr{L}$ is ample…」，未带「$\pi$-」前缀，疑当为「$\pi$-ample」之脱漏（按绝对丰富解读则非真命题），兹照录原书。——译者注}，则 $\pi$ 必然是\trQuasicompact{}的，并且（由\mbox{定理~16.2.6}）是\trSeparated{}的；若 $\pi$ 是仿射的，则一切\trInvertibleSheaf{}都是\trAmple{}的；而若 $\pi$ 是射影的，则相应的 $\OO(1)$ 是\trAmple{}的。由\mbox{练习~16.2.M}，$\mathscr{L}$ 是 $\pi$-\trAmple{}的，当且仅当 $\mathscr{L}^{\otimes n}$ 是 $\pi$-\trAmple{}的；并且若 $Z\hookrightarrow X$ 是一个\trClosedEmbedding，则 $\mathscr{L}|_Z$ 在 $Y$ 上是\trAmple{}的。

由\mbox{定理~16.2.6}(d)，我们有一个自然的\trOpenEmbedding{} $X\hookrightarrow\mathit{Proj}_Y\oplus\pi_*\mathscr{L}^{\otimes\mathrm{d}}$。（你看得出来这个映射是什么吗？另外要当心：$\oplus\pi_*\mathscr{L}^{\otimes\mathrm{d}}$ 未必是一个有限生成的\trSheafOfGradedAlgebras，所以我们这是在通常假设之一不成立的情形下使用 $\mathit{Proj}$ 构造。）

相对整体生成与相对丰富性这些观念，最有用的场合是在\trProper{}的设定中，理由是\mbox{定理~16.2.2}。

\textbf{17.3.6.}\quad \S16.1 的许多陈述可以原封不动地照搬过来。例如，我们有下面这组结果。设 $\pi\colon X\to Y$ 是\trProper{}的，$\mathscr{F}$ 与 $\mathscr{G}$ 是 $X$ 上的\trQCSheaf，而 $\mathscr{L}$ 与 $\mathscr{M}$ 是 $X$ 上的\trInvertibleSheaf。若 $\pi$ 是仿射的，则 $\mathscr{F}$ 是相对\trGloballyGenerated（容易的\mbox{练习~16.1.B}）。若 $\mathscr{F}$ 与 $\mathscr{G}$ 都是相对\trGloballyGenerated，则 $\mathscr{F}\otimes\mathscr{G}$ 亦然（容易的\mbox{练习~16.1.C}）。若 $\mathscr{L}$ 是 $\pi$-极丰富的（\mbox{定义~17.2.2}），则它是 $\pi$-\trBasePointFree{}的（容易的\mbox{练习~16.2.B}）。若 $\mathscr{L}$ 是 $\pi$-极丰富的，而 $\mathscr{M}$ 是 $\pi$-\trBasePointFree{}的（比如说，若它是 $\pi$-极丰富的），则 $\mathscr{L}\otimes\mathscr{M}$ 是 $\pi$-极丰富的（\mbox{练习~16.2.C}）。\mbox{练习~16.2.G} 可以立即推广，以证明：若
\[
\begin{tikzcd}[row sep=2.6em, column sep=4.2em]
X \ar[rr, "\pi"] \ar[dr, "\tau"'] && Y \ar[dl, "\rho"'] \\
& S &
\end{tikzcd}
\]
中的 $\pi$ 是 $S$-概形之间的一个\trFinite\trMorphism，而 $\mathscr{L}$ 是 $Y$ 上一个 $\rho$-\trAmple{}的\trInvertibleSheaf，则 $\pi^*\mathscr{L}$ 是 $\tau$-\trAmple{}的。

出于这些陈述的本性，\S16.1 的某些陈述会需要对 $Y$ 添加拟紧性假设，或者作别的修补。例如：

\textbf{17.3.7.\ \trhTheorem{}。}\ ——\ \emph{设 $\pi\colon X\to Y$ 是\trProper{}的，$\mathscr{L}$ 是 $X$ 上的一个\trInvertibleSheaf，而 $Y$ 是\trQuasicompact{}的。则以下诸条件等价。}

\emph{(a) 对某个 $\mathrm{N}>0$，$\mathscr{L}^{\otimes\mathrm{N}}$ 是 $\pi$-极丰富的。}

\emph{(a') 对一切 $n\gg 0$，$\mathscr{L}^{\otimes n}$ 是 $\pi$-极丰富的。}

\emph{(b) 对一切\trFiniteType{}的\trQCSheaf{} $\mathscr{F}$，存在某个 $n_0$，使得对 $n\ge n_0$，$\mathscr{F}\otimes\mathscr{L}^{\otimes n}$ 都是相对\trGloballyGenerated。}

\emph{(c) \trInvertibleSheaf{} $\mathscr{L}$ 是 $\pi$-\trAmple{}的。}

\textbf{17.3.G.（练习）}\quad 用\mbox{定理~16.2.2} 证明\mbox{定理~17.3.7}。（不重要的评注：\mbox{定理~16.2.2} 所给出的证明用到了诺特假设，但按那里的表述，这些假设是可以去掉的。）

做完上面的练习之后，倘若你需要去掉 $Y$ 上的拟紧性假设，\mbox{定理~17.3.7} 的陈述该如何调整，将会是清楚的。

\textbf{17.3.H.（练习：极丰富性的一个有用的等价定义）}\quad 设 $\pi\colon X\to Y$ 是一个\trProper\trMorphism，$Y$ 是局部\trNoetherian{}的（从而 $X$ 亦然，因为 $\pi$ 是\trFiniteType{}的），而 $\mathscr{L}$ 是 $X$ 上的一个\trInvertibleSheaf。假设你知道在这一情形下 $\pi_*\mathscr{L}$ 是\trFiniteType{}的。（我们日后将如\mbox{练习~17.3.F} 所述证明这一点。诺特假设正是在此处被用到的。）证明：$\mathscr{L}$ 是 $\pi$-极丰富的，当且仅当（i）$\mathscr{L}$ 是相对\trBasePointFree{}的，并且（ii）\mbox{练习~17.3.F} 的那个典范\trMorphism{} $\psi\colon X\to\mathit{Proj}_Y\Sym^\bullet(\pi_*\mathscr{L})$ 是一个\trClosedEmbedding。由此断结：若 $\pi$ 是\trProper{}的，则相对极丰富性这一观念在 $Y$ 上是仿射局部的（可以在\emph{任一份} $Y$ 的仿射覆盖上检验）。\footnote{原书此处作「any affine cover Y」，当为「any affine cover of Y」之脱漏，译文径顺化。——译者注}

\textbf{17.3.8.}\quad 作为推论，\mbox{定理~17.3.7} 蕴含：相对丰富性这一观念在 $Y$ 上是仿射局部的（若 $\pi$ 是\trProper{}的且 $Y$ 是局部\trNoetherian{}的）。

\textbf{17.3.I.（重要而富有挑战性的练习：射影态射的复合通常是射影的）}\quad 设 $\pi\colon X\to Y$ 与 $\rho\colon Y\to Z$ 是两个射影态射，而 $Z$ 是仿射的或\trNoetherian{}的。证明：$\rho\circ\pi$ 是射影的。提示：\mbox{练习~17.3.A}(b) 中给出的射影性判别法将会派上用场。（i）先处理 $Z$ 为仿射的情形。构造如下的\trCommutativeDiagram，从而找出一个在 $Z$ 上的\trClosedEmbedding{} $X\hookrightarrow\mathbb{P}\mathscr{F}^{\oplus n}$。（$\mathbb{P}\mathscr{V}$ 的定义见 \S17.2.3，其中还包括一处关于对偶的警示。）在这个图中，所有的包含都是\trClosedEmbedding，而所有的花体字母指的都是\trFiniteType{}的\trQCSheaf。
\[
\begin{tikzcd}[row sep=3.1em, column sep=1.7em]
X \ar[r, hook] \ar[drr, "\pi"'] & \mathbb{P}\mathscr{G} \ar[r, hook, "(\dagger)"] \ar[dr] & \mathbb{P}^{n-1}_Z\times_Z Y \ar[r, hook] \ar[d] & \mathbb{P}^{n-1}_Z\times_Z\mathbb{P}\mathscr{F} \ar[r, hook, "\mathrm{Segre}", "\text{参\mbox{练习~10.6.C、10.6.D}}"'] \ar[d] & \mathbb{P}(\mathscr{F}^{\oplus n}) \ar[ddll] \\
&& Y \ar[r, hook] \ar[d, "\rho"'] & \mathbb{P}\mathscr{F} \ar[dl] & \\
&& Z & &
\end{tikzcd}
\]
按如下方式构造这个\trClosedEmbedding{} $(\dagger)$。设 $\mathscr{M}$ 是 $Y$ 在 $Z$ 上的那个极丰富线丛。于是 $\mathscr{M}$ 是\trAmple{}的，从而由\mbox{定理~16.2.2}，对 $\mathrm{m}\gg 0$，$\mathscr{G}\otimes\mathscr{M}^{\otimes\mathrm{m}}$ 由有限多个\trGlobalSections{}生成。设 $\OO_Y^{\oplus n}\twoheadrightarrow\mathscr{G}\otimes\mathscr{M}^{\otimes\mathrm{m}}$ 是相应的满射。这就诱导出一个\trClosedEmbedding{} $\mathbb{P}(\mathscr{G}\otimes\mathscr{M}^{\otimes\mathrm{m}})\hookrightarrow\mathbb{P}^{n-1}_Y$。但 $\mathbb{P}(\mathscr{G}\otimes\mathscr{M}^{\otimes\mathrm{m}})\cong\mathbb{P}\mathscr{G}$（\mbox{练习~17.2.G}），而 $\mathbb{P}^{n-1}_Y=\mathbb{P}^{n-1}_Z\times_Z Y$。（ii）展开这个图以证明：（对 $Z$ 仿射的情形）若 $\mathscr{L}$ 是 $\pi$-极丰富的且 $\mathscr{M}$ 是 $\rho$-极丰富的，则对 $\mathrm{m}\gg 0$，$\mathscr{L}\otimes\pi^*\mathscr{M}^{\otimes\mathrm{m}}$ 是 $(\rho\circ\pi)$-极丰富的。然后，借助\mbox{练习~17.3.H}，通过用有限多个仿射开子集覆盖 $Z$ 来处理一般情形。

\textbf{17.3.9.}\ \extrahard\ \textbf{\trAmple\trVectorBundle。}\quad \textbf{\trAmple\trVectorBundle}（ample vector bundle）的观念在文献的某些部分中是有用的，所以我们在此给出它的定义，尽管我们不会用到这个观念。设 $\mathscr{E}$ 是\trProper{} $A$-\trScheme{} $X$ 上的一个\trLocallyFreeSheaf。若 $\OO_{\mathit{Proj}_X\Sym^\bullet\mathscr{E}/X}(1)$ 是一个\trAmple{}的\trInvertibleSheaf，则称 $\mathscr{E}$ 是\textbf{\trAmple{}的}（ample）。特别地，利用\mbox{练习~17.2.G}，你可以验证：一个\trInvertibleSheaf{}作为\trLocallyFreeSheaf（按这个定义）是\trAmple{}的，当且仅当它作为\trInvertibleSheaf（\mbox{定义~16.2.1}）是\trAmple{}的，从而免于一场记号危机。（这个\trProper{}假设是可以放宽的；之所以把它收进来，只是因为 \mbox{定义~16.2.1} 的丰富性仅对\trProper{}概形而言。）

\textbf{17.3.10.}\ \extrahard\ \textbf{拟射影态射。}\quad 与射影及拟射影 $A$-\trScheme（\S4.5.10）相类比，我们可以定义拟射影态射。\emph{若 $Y$ 是\trQuasicompact{}的}，我们就说 $\pi\colon X\to Y$ 是一个\textbf{拟射影态射}（quasiprojective morphism），若 $\pi$ 可以表达成一个到某个在 $Y$ 上射影的概形的\trQuasicompact\trOpenEmbedding。这并不是一个多么好的观念，我们也不会用到它。（拟射影态射的一般定义略显微妙——见 [Gr-EGA, II.5.3]——而我们用不着它。）

\textbf{17.3.11.}\ \extrahard\ \textbf{\trQuasiaffine{}态射。}

既然我们已经引入了拟射影态射，就不妨顺带简短地引入\trQuasiaffine{}态射（以及\trQuasiaffine{}概形），因为某些读者或许有用到它们的场合。这些观念中有许多本可以早在很久以前就引入，但由于我们永远不会用到它们，我们索性在此一并处理。

一个\trScheme{} $X$ 是\textbf{\trQuasiaffine{}的}（quasiaffine），若它容许一个到某个\trAffineScheme{}的\trQuasicompact\trOpenEmbedding。这蕴含 $X$ 是\trQuasicompact{}且\trSeparated{}的。若 $X$ 是\trNoetherian{}的（对多数人而言这是最要紧的情形），则任何\trOpenEmbedding{}当然是自动\trQuasicompact{}的。

\textbf{17.3.J.（练习）}\quad 证明：$X$ 是\trQuasiaffine{}的，当且仅当典范映射 $X\to\Spec\Gamma(X,\OO_X)$（定义见\mbox{练习~7.3.G} 及其后那段话）是一个\trQuasicompact\trOpenEmbedding。于是，\trQuasiaffine{}概形随之自带一个\emph{典范的}到\trAffineScheme{}的\trQuasicompact\trOpenEmbedding。提示：为方便计，令 $A=\Gamma(X,\OO_X)$。设 $X\hookrightarrow\Spec R$ 是一个\trQuasicompact\trOpenEmbedding。我们希望证明 $X\to\Spec A$ 是一个\trQuasicompact\trOpenEmbedding。把 $X\hookrightarrow\Spec R$ 经由 $X\to\Spec A\to\Spec R$ 来分解。对任意选定的点 $p\in X$，证明 $X\to\Spec A$ 在其某个\trOpen\trNeighborhood{}中是一个开嵌入，做法如下。选 $r\in R$ 使得 $p\in\mathrm{D}(r)\subset X$。注意：若 $X_r=\{q\in X: r(q)\neq 0\}$，则 $\Gamma(X_r,\OO_X)=\Gamma(X,\OO_X)_r$（由\mbox{练习~15.4.N}，这里用到了 $X$ 是\trQuasicompact{}且\trQuasiSeparated{}的这一事实）。用它来证明映射 $X_r\to\Spec A_r$ 是一个\trIsomorphism。

不难证明：\emph{$X$} 是\trQuasiaffine{}的，当且仅当 $\OO_X$ \emph{是}\trAmple{}的，但我们不会用到这一事实。

一个\trMorphism{} $\pi\colon X\to Y$ 称为\textbf{\trQuasiaffine{}的}，若 $Y$ 的每个\trAffine{}开子集的\trInverseImage{}都是一个\trQuasiaffine{}概形。由\mbox{练习~17.3.J}，这等价于：$\pi$ 是\trQuasicompact{}且\trSeparated{}的，并且自然映射 $X\to\mathit{Spec}\,\pi_*\OO_X$ 是一个\trQuasicompact\trOpenEmbedding。这蕴含拟仿射性（quasiaffineness）这一观念在靶上是局部的（可以在一个\trOpenCover{}上检验），并且在靶上还是仿射局部的（可以选取一个仿射覆盖，然后检验这些开集的原像都是\trQuasiaffine{}的）。\trQuasiaffine{}态射在\trBaseChange{}下保持：若 $Y$ 上的一个\trMorphism{} $X\hookrightarrow Z$ 是到某个在 $Y$ 上仿射的概形的\trQuasicompact\trOpenEmbedding，则对任意 $W\to Y$，$X\times_Y W\hookrightarrow Z\times_Y W$ 是到某个在 $W$ 上仿射的概形的\trQuasicompact\trOpenEmbedding。（有趣的是，\mbox{练习~17.3.J} \emph{并}不是用来证明这一\trBaseChange{}性质的正确工具。）

容易验证，\trQuasiaffine{}态射在复合下保持 [Stacks, tag 01SN]（于是\trQuasiaffine{}态射构成一个「\trReasonable」类）。从而，\trQuasicompact{}的\trLocallyClosedEmbedding{}都是\trQuasiaffine{}的。若 $X$ 是仿射的，则 $X\to Y$ 是\trQuasiaffine{}的，当且仅当它是\trQuasicompact{}的（因为 $Y$ 的任何仿射开子集的原像都是仿射概形 $X$ 的一个开子集）。特别地，由针对\trQuasicompact{}态射的\trCancellationTheorem~11.1.1，从仿射概形出发到拟分离概形的任何\trMorphism{}都是\trQuasiaffine{}的。

